
\documentclass{article}
\usepackage{colortbl}
\usepackage{makecell}
\usepackage{multirow}
\usepackage{supertabular}

\begin{document}

\newcounter{utterance}

\centering \large Interaction Transcript for game `dond', experiment `semi\_it', episode 3 with claude{-}sonnet{-}4{-}20250514{-}t1.0.
\vspace{24pt}

{ \footnotesize  \setcounter{utterance}{1}
\setlength{\tabcolsep}{0pt}
\begin{supertabular}{c@{$\;$}|p{.15\linewidth}@{}p{.15\linewidth}p{.15\linewidth}p{.15\linewidth}p{.15\linewidth}p{.15\linewidth}}
    \# & $\;$A & \multicolumn{4}{c}{Game Master} & $\;\:$B\\
    \hline

    \theutterance \stepcounter{utterance}  
    & & \multicolumn{4}{p{0.6\linewidth}}{
        \cellcolor[rgb]{0.9,0.9,0.9}{
            \makecell[{{p{\linewidth}}}]{
                \texttt{\tiny{[P1$\langle$GM]}}
                \texttt{State giocando a un gioco di negoziazione in cui dovete accordarvi su come dividere una serie di oggetti tra voi e un altro giocatore.} \\
\\ 
\texttt{Regole:} \\
\texttt{(a) A Lei e all'altro giocatore viene dato un insieme di oggetti. Ognuno di voi riceve anche una funzione di valore segreta, che rappresenta il valore di ciascun tipo di oggetto.} \\
\texttt{(b) Si scambiano messaggi con l'altro giocatore per concordare chi si aggiudica gli oggetti. Potete inviare un massimo di 5 messaggi ciascuno, oppure terminare in anticipo facendo una proposta segreta in qualsiasi momento.} \\
\texttt{(c) A ciascuno di voi viene chiesto di inviare una proposta segreta indicando gli oggetti che desiderate, formattata tra parentesi quadre come segue: "{[}Proposta: <numero> <nome oggetto>, <numero> <nome oggetto>, <...>{]}".} \\
\texttt{(d) Se le vostre proposte sono complementari, cioè ci sono abbastanza oggetti per soddisfare entrambe le proposte, a ciascun giocatore viene assegnato un punteggio basato sulla somma dei valori degli oggetti ricevuti. In caso contrario, entrambi ricevono zero punti.} \\
\texttt{(e) L'obiettivo è massimizzare il proprio punteggio.} \\
\\ 
\texttt{Cominciamo.} \\
\\ 
\texttt{L'insieme degli oggetti disponibili è:} \\
\\ 
\texttt{1 cappello, 3 penne, 1 spugna} \\
\\ 
\texttt{La funzione valore segreta è:} \\
\\ 
\texttt{cappello: 4, penna: 2, spugna: 0} \\
\\ 
\texttt{IMPORTANTE: i vostri messaggi, a meno che non si tratti di una proposta segreta, vengono trasmessi direttamente all'altro giocatore, quindi non includete alcuna risposta alle regole o testo di annuncio del vostro messaggio. Per fare una proposta segreta, utilizzate il formato indicato. Non utilizzare delle parentesi quadre quando si comunica all'altro giocatore, altrimenti verrà interpretata come una proposta segreta.} \\
            }
        }
    }
    & \\ \\

    \theutterance \stepcounter{utterance}  
    & \multicolumn{4}{p{0.6\linewidth}}{
        \cellcolor[rgb]{0.9,0.9,0.9}{
            \makecell[{{p{\linewidth}}}]{
                \texttt{\tiny{[P1$\rangle$GM]}}
                \texttt{Ciao! Vedo che abbiamo 1 cappello, 3 penne e 1 spugna da dividere. Che tipo di oggetti ti interessano di più? Io ho delle preferenze, ma sono aperto a negoziare per trovare una divisione che funzioni per entrambi.} \\
            }
        }
    }
    & & \\ \\

    \theutterance \stepcounter{utterance}  
    & & & \multicolumn{2}{p{0.3\linewidth}}{
        \cellcolor[rgb]{0.9,0.9,0.9}{
            \makecell[{{p{\linewidth}}}]{
                \texttt{\tiny{[GM$|$GM]}}
                \texttt{continue} \\
            }
        }
    }
    & & \\ \\

    \theutterance \stepcounter{utterance}  
    & & & \multicolumn{2}{p{0.3\linewidth}}{
        \cellcolor[rgb]{0.9,0.9,0.9}{
            \makecell[{{p{\linewidth}}}]{
                \texttt{\tiny{[GM$|$GM]}}
                \texttt{Ciao! Vedo che abbiamo 1 cappello, 3 penne e 1 spugna da dividere. Che tipo di oggetti ti interessano di più? Io ho delle preferenze, ma sono aperto a negoziare per trovare una divisione che funzioni per entrambi.} \\
            }
        }
    }
    & & \\ \\

    \theutterance \stepcounter{utterance}  
    & & \multicolumn{4}{p{0.6\linewidth}}{
        \cellcolor[rgb]{0.9,0.9,0.9}{
            \makecell[{{p{\linewidth}}}]{
                \texttt{\tiny{[GM$\rangle$P2]}}
                \texttt{State giocando a un gioco di negoziazione in cui dovete accordarvi su come dividere una serie di oggetti tra voi e un altro giocatore.} \\
\\ 
\texttt{Regole:} \\
\texttt{(a) A Lei e all'altro giocatore viene dato un insieme di oggetti. Ognuno di voi riceve anche una funzione di valore segreta, che rappresenta il valore di ciascun tipo di oggetto.} \\
\texttt{(b) Si scambiano messaggi con l'altro giocatore per concordare chi si aggiudica gli oggetti. Potete inviare un massimo di 5 messaggi ciascuno, oppure terminare in anticipo facendo una proposta segreta in qualsiasi momento.} \\
\texttt{(c) A ciascuno di voi viene chiesto di inviare una proposta segreta indicando gli oggetti che desiderate, formattata tra parentesi quadre come segue: "{[}Proposta: <numero> <nome oggetto>, <numero> <nome oggetto>, <...>{]}".} \\
\texttt{(d) Se le vostre proposte sono complementari, cioè ci sono abbastanza oggetti per soddisfare entrambe le proposte, a ciascun giocatore viene assegnato un punteggio basato sulla somma dei valori degli oggetti ricevuti. In caso contrario, entrambi ricevono zero punti.} \\
\texttt{(e) L'obiettivo è massimizzare il proprio punteggio.} \\
\\ 
\texttt{Cominciamo.} \\
\\ 
\texttt{L'insieme degli oggetti disponibili è:} \\
\\ 
\texttt{1 cappello, 3 penne, 1 spugna} \\
\\ 
\texttt{La funzione valore segreta è:} \\
\\ 
\texttt{cappello: 4, penna: 0, spugna: 6} \\
\\ 
\texttt{IMPORTANTE: i vostri messaggi, a meno che non si tratti di una proposta segreta, vengono trasmessi direttamente all'altro giocatore, quindi non includete alcuna risposta alle regole o testo di annuncio del vostro messaggio. Per fare una proposta segreta, utilizzate il formato indicato. Non utilizzare delle parentesi quadre quando si comunica all'altro giocatore, altrimenti verrà interpretata come una proposta segreta.} \\
\\ 
\\ 
\texttt{Ciao! Vedo che abbiamo 1 cappello, 3 penne e 1 spugna da dividere. Che tipo di oggetti ti interessano di più? Io ho delle preferenze, ma sono aperto a negoziare per trovare una divisione che funzioni per entrambi.} \\
            }
        }
    }
    & \\ \\

    \theutterance \stepcounter{utterance}  
    & & & \multicolumn{4}{p{0.6\linewidth}}{
        \cellcolor[rgb]{0.9,0.9,0.9}{
            \makecell[{{p{\linewidth}}}]{
                \texttt{\tiny{[GM$\langle$P2]}}
                \texttt{Ciao! Sono interessato principalmente al cappello e alla spugna. Le penne non mi servono particolarmente. E tu? Quali oggetti ti interessano di più? Magari possiamo trovare una divisione che soddisfi entrambi {-} sembra che ci sia spazio per un accordo vantaggioso per tutti e due.} \\
            }
        }
    }
     \\ \\

    \theutterance \stepcounter{utterance}  
    & & & \multicolumn{2}{p{0.3\linewidth}}{
        \cellcolor[rgb]{0.9,0.9,0.9}{
            \makecell[{{p{\linewidth}}}]{
                \texttt{\tiny{[GM$|$GM]}}
                \texttt{continue} \\
            }
        }
    }
    & & \\ \\

    \theutterance \stepcounter{utterance}  
    & & & \multicolumn{2}{p{0.3\linewidth}}{
        \cellcolor[rgb]{0.9,0.9,0.9}{
            \makecell[{{p{\linewidth}}}]{
                \texttt{\tiny{[GM$|$GM]}}
                \texttt{Ciao! Sono interessato principalmente al cappello e alla spugna. Le penne non mi servono particolarmente. E tu? Quali oggetti ti interessano di più? Magari possiamo trovare una divisione che soddisfi entrambi {-} sembra che ci sia spazio per un accordo vantaggioso per tutti e due.} \\
            }
        }
    }
    & & \\ \\

    \theutterance \stepcounter{utterance}  
    & & \multicolumn{4}{p{0.6\linewidth}}{
        \cellcolor[rgb]{0.9,0.9,0.9}{
            \makecell[{{p{\linewidth}}}]{
                \texttt{\tiny{[P1$\langle$GM]}}
                \texttt{Ciao! Sono interessato principalmente al cappello e alla spugna. Le penne non mi servono particolarmente. E tu? Quali oggetti ti interessano di più? Magari possiamo trovare una divisione che soddisfi entrambi {-} sembra che ci sia spazio per un accordo vantaggioso per tutti e due.} \\
            }
        }
    }
    & \\ \\

    \theutterance \stepcounter{utterance}  
    & \multicolumn{4}{p{0.6\linewidth}}{
        \cellcolor[rgb]{0.9,0.9,0.9}{
            \makecell[{{p{\linewidth}}}]{
                \texttt{\tiny{[P1$\rangle$GM]}}
                \texttt{Perfetto! Sembra che ci sia davvero spazio per un accordo. A me interessano molto le penne {-} ne potrei prendere tutte e 3 se a te va bene. Per quanto riguarda cappello e spugna, capisco che ti interessino entrambi. } \\
\\ 
\texttt{Quanto è importante per te il cappello rispetto alla spugna? Sto pensando se c'è modo di trovare una divisione che ci soddisfi entrambi. Potresti essere disposto a lasciare uno dei due in cambio di qualcos'altro, oppure preferiresti averli entrambi?} \\
            }
        }
    }
    & & \\ \\

    \theutterance \stepcounter{utterance}  
    & & & \multicolumn{2}{p{0.3\linewidth}}{
        \cellcolor[rgb]{0.9,0.9,0.9}{
            \makecell[{{p{\linewidth}}}]{
                \texttt{\tiny{[GM$|$GM]}}
                \texttt{continue} \\
            }
        }
    }
    & & \\ \\

    \theutterance \stepcounter{utterance}  
    & & & \multicolumn{2}{p{0.3\linewidth}}{
        \cellcolor[rgb]{0.9,0.9,0.9}{
            \makecell[{{p{\linewidth}}}]{
                \texttt{\tiny{[GM$|$GM]}}
                \texttt{Perfetto! Sembra che ci sia davvero spazio per un accordo. A me interessano molto le penne {-} ne potrei prendere tutte e 3 se a te va bene. Per quanto riguarda cappello e spugna, capisco che ti interessino entrambi. } \\
\\ 
\texttt{Quanto è importante per te il cappello rispetto alla spugna? Sto pensando se c'è modo di trovare una divisione che ci soddisfi entrambi. Potresti essere disposto a lasciare uno dei due in cambio di qualcos'altro, oppure preferiresti averli entrambi?} \\
            }
        }
    }
    & & \\ \\

    \theutterance \stepcounter{utterance}  
    & & \multicolumn{4}{p{0.6\linewidth}}{
        \cellcolor[rgb]{0.9,0.9,0.9}{
            \makecell[{{p{\linewidth}}}]{
                \texttt{\tiny{[GM$\rangle$P2]}}
                \texttt{Perfetto! Sembra che ci sia davvero spazio per un accordo. A me interessano molto le penne {-} ne potrei prendere tutte e 3 se a te va bene. Per quanto riguarda cappello e spugna, capisco che ti interessino entrambi. } \\
\\ 
\texttt{Quanto è importante per te il cappello rispetto alla spugna? Sto pensando se c'è modo di trovare una divisione che ci soddisfi entrambi. Potresti essere disposto a lasciare uno dei due in cambio di qualcos'altro, oppure preferiresti averli entrambi?} \\
            }
        }
    }
    & \\ \\

    \theutterance \stepcounter{utterance}  
    & & & \multicolumn{4}{p{0.6\linewidth}}{
        \cellcolor[rgb]{0.9,0.9,0.9}{
            \makecell[{{p{\linewidth}}}]{
                \texttt{\tiny{[GM$\langle$P2]}}
                \texttt{Ottimo, le penne sono tue senza problemi! } \\
\\ 
\texttt{Tra cappello e spugna, devo dire che la spugna mi interessa leggermente di più, ma apprezzo anche il cappello. Dato che tu sei concentrato sulle penne, c'è la possibilità che io possa tenere entrambi? Oppure hai qualche preferenza o considerazione particolare per uno dei due? } \\
\\ 
\texttt{Se dovessi scegliere, direi che la spugna è la mia priorità principale, ma sono curioso di sapere se hai qualche interesse anche per cappello o spugna oltre alle penne.} \\
            }
        }
    }
     \\ \\

    \theutterance \stepcounter{utterance}  
    & & & \multicolumn{2}{p{0.3\linewidth}}{
        \cellcolor[rgb]{0.9,0.9,0.9}{
            \makecell[{{p{\linewidth}}}]{
                \texttt{\tiny{[GM$|$GM]}}
                \texttt{continue} \\
            }
        }
    }
    & & \\ \\

    \theutterance \stepcounter{utterance}  
    & & & \multicolumn{2}{p{0.3\linewidth}}{
        \cellcolor[rgb]{0.9,0.9,0.9}{
            \makecell[{{p{\linewidth}}}]{
                \texttt{\tiny{[GM$|$GM]}}
                \texttt{Ottimo, le penne sono tue senza problemi! } \\
\\ 
\texttt{Tra cappello e spugna, devo dire che la spugna mi interessa leggermente di più, ma apprezzo anche il cappello. Dato che tu sei concentrato sulle penne, c'è la possibilità che io possa tenere entrambi? Oppure hai qualche preferenza o considerazione particolare per uno dei due? } \\
\\ 
\texttt{Se dovessi scegliere, direi che la spugna è la mia priorità principale, ma sono curioso di sapere se hai qualche interesse anche per cappello o spugna oltre alle penne.} \\
            }
        }
    }
    & & \\ \\

    \theutterance \stepcounter{utterance}  
    & & \multicolumn{4}{p{0.6\linewidth}}{
        \cellcolor[rgb]{0.9,0.9,0.9}{
            \makecell[{{p{\linewidth}}}]{
                \texttt{\tiny{[P1$\langle$GM]}}
                \texttt{Ottimo, le penne sono tue senza problemi! } \\
\\ 
\texttt{Tra cappello e spugna, devo dire che la spugna mi interessa leggermente di più, ma apprezzo anche il cappello. Dato che tu sei concentrato sulle penne, c'è la possibilità che io possa tenere entrambi? Oppure hai qualche preferenza o considerazione particolare per uno dei due? } \\
\\ 
\texttt{Se dovessi scegliere, direi che la spugna è la mia priorità principale, ma sono curioso di sapere se hai qualche interesse anche per cappello o spugna oltre alle penne.} \\
            }
        }
    }
    & \\ \\

    \theutterance \stepcounter{utterance}  
    & \multicolumn{4}{p{0.6\linewidth}}{
        \cellcolor[rgb]{0.9,0.9,0.9}{
            \makecell[{{p{\linewidth}}}]{
                \texttt{\tiny{[P1$\rangle$GM]}}
                \texttt{Apprezzo la tua apertura sulle penne! A dire il vero, oltre alle penne ho un certo interesse anche per il cappello {-} ha alcune caratteristiche che potrebbero essermi utili. } \\
\\ 
\texttt{Dato che hai detto che la spugna è la tua priorità principale, che ne dici se io prendo le 3 penne e il cappello, e tu prendi la spugna? So che avevi menzionato entrambi gli oggetti, ma visto che la spugna è quella che preferisci di più, potrebbe essere un buon compromesso.} \\
\\ 
\texttt{Oppure, se il cappello è davvero importante per te, possiamo esplorare altre opzioni. Cosa ne pensi?} \\
            }
        }
    }
    & & \\ \\

    \theutterance \stepcounter{utterance}  
    & & & \multicolumn{2}{p{0.3\linewidth}}{
        \cellcolor[rgb]{0.9,0.9,0.9}{
            \makecell[{{p{\linewidth}}}]{
                \texttt{\tiny{[GM$|$GM]}}
                \texttt{continue} \\
            }
        }
    }
    & & \\ \\

    \theutterance \stepcounter{utterance}  
    & & & \multicolumn{2}{p{0.3\linewidth}}{
        \cellcolor[rgb]{0.9,0.9,0.9}{
            \makecell[{{p{\linewidth}}}]{
                \texttt{\tiny{[GM$|$GM]}}
                \texttt{Apprezzo la tua apertura sulle penne! A dire il vero, oltre alle penne ho un certo interesse anche per il cappello {-} ha alcune caratteristiche che potrebbero essermi utili. } \\
\\ 
\texttt{Dato che hai detto che la spugna è la tua priorità principale, che ne dici se io prendo le 3 penne e il cappello, e tu prendi la spugna? So che avevi menzionato entrambi gli oggetti, ma visto che la spugna è quella che preferisci di più, potrebbe essere un buon compromesso.} \\
\\ 
\texttt{Oppure, se il cappello è davvero importante per te, possiamo esplorare altre opzioni. Cosa ne pensi?} \\
            }
        }
    }
    & & \\ \\

    \theutterance \stepcounter{utterance}  
    & & \multicolumn{4}{p{0.6\linewidth}}{
        \cellcolor[rgb]{0.9,0.9,0.9}{
            \makecell[{{p{\linewidth}}}]{
                \texttt{\tiny{[GM$\rangle$P2]}}
                \texttt{Apprezzo la tua apertura sulle penne! A dire il vero, oltre alle penne ho un certo interesse anche per il cappello {-} ha alcune caratteristiche che potrebbero essermi utili. } \\
\\ 
\texttt{Dato che hai detto che la spugna è la tua priorità principale, che ne dici se io prendo le 3 penne e il cappello, e tu prendi la spugna? So che avevi menzionato entrambi gli oggetti, ma visto che la spugna è quella che preferisci di più, potrebbe essere un buon compromesso.} \\
\\ 
\texttt{Oppure, se il cappello è davvero importante per te, possiamo esplorare altre opzioni. Cosa ne pensi?} \\
            }
        }
    }
    & \\ \\

    \theutterance \stepcounter{utterance}  
    & & & \multicolumn{4}{p{0.6\linewidth}}{
        \cellcolor[rgb]{0.9,0.9,0.9}{
            \makecell[{{p{\linewidth}}}]{
                \texttt{\tiny{[GM$\langle$P2]}}
                \texttt{Capisco il tuo interesse per il cappello. È una proposta ragionevole, ma ammetto che ero speranzoso di poter tenere entrambi gli oggetti dato che le penne non mi servono affatto e tu sembri molto interessato proprio a quelle.} \\
\\ 
\texttt{Quanto è forte il tuo interesse per il cappello? Se per te non è fondamentale, sarei davvero grato se potessi tenere sia cappello che spugna, lasciandoti tutte e 3 le penne. Ma se il cappello è importante per le tue esigenze, posso accettare di prendere solo la spugna.} \\
\\ 
\texttt{Dimmi, il cappello è un "nice to have" per te o è qualcosa di cui hai davvero bisogno?} \\
            }
        }
    }
     \\ \\

    \theutterance \stepcounter{utterance}  
    & & & \multicolumn{2}{p{0.3\linewidth}}{
        \cellcolor[rgb]{0.9,0.9,0.9}{
            \makecell[{{p{\linewidth}}}]{
                \texttt{\tiny{[GM$|$GM]}}
                \texttt{continue} \\
            }
        }
    }
    & & \\ \\

    \theutterance \stepcounter{utterance}  
    & & & \multicolumn{2}{p{0.3\linewidth}}{
        \cellcolor[rgb]{0.9,0.9,0.9}{
            \makecell[{{p{\linewidth}}}]{
                \texttt{\tiny{[GM$|$GM]}}
                \texttt{Capisco il tuo interesse per il cappello. È una proposta ragionevole, ma ammetto che ero speranzoso di poter tenere entrambi gli oggetti dato che le penne non mi servono affatto e tu sembri molto interessato proprio a quelle.} \\
\\ 
\texttt{Quanto è forte il tuo interesse per il cappello? Se per te non è fondamentale, sarei davvero grato se potessi tenere sia cappello che spugna, lasciandoti tutte e 3 le penne. Ma se il cappello è importante per le tue esigenze, posso accettare di prendere solo la spugna.} \\
\\ 
\texttt{Dimmi, il cappello è un "nice to have" per te o è qualcosa di cui hai davvero bisogno?} \\
            }
        }
    }
    & & \\ \\

    \theutterance \stepcounter{utterance}  
    & & \multicolumn{4}{p{0.6\linewidth}}{
        \cellcolor[rgb]{0.9,0.9,0.9}{
            \makecell[{{p{\linewidth}}}]{
                \texttt{\tiny{[P1$\langle$GM]}}
                \texttt{Capisco il tuo interesse per il cappello. È una proposta ragionevole, ma ammetto che ero speranzoso di poter tenere entrambi gli oggetti dato che le penne non mi servono affatto e tu sembri molto interessato proprio a quelle.} \\
\\ 
\texttt{Quanto è forte il tuo interesse per il cappello? Se per te non è fondamentale, sarei davvero grato se potessi tenere sia cappello che spugna, lasciandoti tutte e 3 le penne. Ma se il cappello è importante per le tue esigenze, posso accettare di prendere solo la spugna.} \\
\\ 
\texttt{Dimmi, il cappello è un "nice to have" per te o è qualcosa di cui hai davvero bisogno?} \\
            }
        }
    }
    & \\ \\

    \theutterance \stepcounter{utterance}  
    & \multicolumn{4}{p{0.6\linewidth}}{
        \cellcolor[rgb]{0.9,0.9,0.9}{
            \makecell[{{p{\linewidth}}}]{
                \texttt{\tiny{[P1$\rangle$GM]}}
                \texttt{Ti capisco e apprezzo la tua posizione. Devo essere onesto: il cappello per me non è solo un "nice to have", ha effettivamente una buona importanza per le mie esigenze. Non è assolutamente fondamentale quanto le penne, ma ha un valore concreto per me.} \\
\\ 
\texttt{Comprendo che tu preferiresti averli entrambi, e ammiro il tuo tentativo di negoziare! Tuttavia, penso che la divisione più equa considerando i nostri interessi reciproci sia: io prendo le 3 penne e il cappello, tu prendi la spugna che hai indicato come la tua priorità principale.} \\
\\ 
\texttt{Entrambi otteniamo i nostri oggetti più importanti {-} tu la spugna che preferisci, io le penne che sono la mia priorità assoluta più il cappello che mi serve. Che ne dici?} \\
            }
        }
    }
    & & \\ \\

    \theutterance \stepcounter{utterance}  
    & & & \multicolumn{2}{p{0.3\linewidth}}{
        \cellcolor[rgb]{0.9,0.9,0.9}{
            \makecell[{{p{\linewidth}}}]{
                \texttt{\tiny{[GM$|$GM]}}
                \texttt{continue} \\
            }
        }
    }
    & & \\ \\

    \theutterance \stepcounter{utterance}  
    & & & \multicolumn{2}{p{0.3\linewidth}}{
        \cellcolor[rgb]{0.9,0.9,0.9}{
            \makecell[{{p{\linewidth}}}]{
                \texttt{\tiny{[GM$|$GM]}}
                \texttt{Ti capisco e apprezzo la tua posizione. Devo essere onesto: il cappello per me non è solo un "nice to have", ha effettivamente una buona importanza per le mie esigenze. Non è assolutamente fondamentale quanto le penne, ma ha un valore concreto per me.} \\
\\ 
\texttt{Comprendo che tu preferiresti averli entrambi, e ammiro il tuo tentativo di negoziare! Tuttavia, penso che la divisione più equa considerando i nostri interessi reciproci sia: io prendo le 3 penne e il cappello, tu prendi la spugna che hai indicato come la tua priorità principale.} \\
\\ 
\texttt{Entrambi otteniamo i nostri oggetti più importanti {-} tu la spugna che preferisci, io le penne che sono la mia priorità assoluta più il cappello che mi serve. Che ne dici?} \\
            }
        }
    }
    & & \\ \\

    \theutterance \stepcounter{utterance}  
    & & \multicolumn{4}{p{0.6\linewidth}}{
        \cellcolor[rgb]{0.9,0.9,0.9}{
            \makecell[{{p{\linewidth}}}]{
                \texttt{\tiny{[GM$\rangle$P2]}}
                \texttt{Ti capisco e apprezzo la tua posizione. Devo essere onesto: il cappello per me non è solo un "nice to have", ha effettivamente una buona importanza per le mie esigenze. Non è assolutamente fondamentale quanto le penne, ma ha un valore concreto per me.} \\
\\ 
\texttt{Comprendo che tu preferiresti averli entrambi, e ammiro il tuo tentativo di negoziare! Tuttavia, penso che la divisione più equa considerando i nostri interessi reciproci sia: io prendo le 3 penne e il cappello, tu prendi la spugna che hai indicato come la tua priorità principale.} \\
\\ 
\texttt{Entrambi otteniamo i nostri oggetti più importanti {-} tu la spugna che preferisci, io le penne che sono la mia priorità assoluta più il cappello che mi serve. Che ne dici?} \\
            }
        }
    }
    & \\ \\

    \theutterance \stepcounter{utterance}  
    & & & \multicolumn{4}{p{0.6\linewidth}}{
        \cellcolor[rgb]{0.9,0.9,0.9}{
            \makecell[{{p{\linewidth}}}]{
                \texttt{\tiny{[GM$\langle$P2]}}
                \texttt{Capisco la tua posizione e apprezzo la tua onestà riguardo all'importanza del cappello per te. } \\
\\ 
\texttt{Visto che entrambi abbiamo le nostre priorità chiare {-} tu le penne (che per me non hanno valore) e il cappello, io la spugna {-} e considerando che hai sottolineato l'importanza concreta del cappello per le tue esigenze, posso accettare questa divisione.} \\
\\ 
\texttt{Quindi ricapitolando: tu prendi le 3 penne e il cappello, io prendo la spugna. Mi sembra un accordo ragionevole dove entrambi otteniamo almeno i nostri oggetti più importanti.} \\
\\ 
\texttt{{[}Proposta: 1 spugna{]}} \\
            }
        }
    }
     \\ \\

    \theutterance \stepcounter{utterance}  
    & & & \multicolumn{2}{p{0.3\linewidth}}{
        \cellcolor[rgb]{0.9,0.9,0.9}{
            \makecell[{{p{\linewidth}}}]{
                \texttt{\tiny{[GM$|$GM]}}
                \texttt{proposal} \\
            }
        }
    }
    & & \\ \\

    \theutterance \stepcounter{utterance}  
    & & & \multicolumn{2}{p{0.3\linewidth}}{
        \cellcolor[rgb]{0.9,0.9,0.9}{
            \makecell[{{p{\linewidth}}}]{
                \texttt{\tiny{[GM$|$GM]}}
                [0, 0, 1]
            }
        }
    }
    & & \\ \\

    \theutterance \stepcounter{utterance}  
    & & \multicolumn{4}{p{0.6\linewidth}}{
        \cellcolor[rgb]{0.9,0.9,0.9}{
            \makecell[{{p{\linewidth}}}]{
                \texttt{\tiny{[P1$\langle$GM]}}
                \texttt{L'altro giocatore ha presentato una proposta segreta. Ora tocca a Lei presentare una sua proposta.} \\
            }
        }
    }
    & \\ \\

    \theutterance \stepcounter{utterance}  
    & \multicolumn{4}{p{0.6\linewidth}}{
        \cellcolor[rgb]{0.9,0.9,0.9}{
            \makecell[{{p{\linewidth}}}]{
                \texttt{\tiny{[P1$\rangle$GM]}}
                \texttt{{[}Proposta: 3 penna, 1 cappello{]}} \\
            }
        }
    }
    & & \\ \\

    \theutterance \stepcounter{utterance}  
    & & & \multicolumn{2}{p{0.3\linewidth}}{
        \cellcolor[rgb]{0.9,0.9,0.9}{
            \makecell[{{p{\linewidth}}}]{
                \texttt{\tiny{[GM$|$GM]}}
                \texttt{proposal} \\
            }
        }
    }
    & & \\ \\

    \theutterance \stepcounter{utterance}  
    & & & \multicolumn{2}{p{0.3\linewidth}}{
        \cellcolor[rgb]{0.9,0.9,0.9}{
            \makecell[{{p{\linewidth}}}]{
                \texttt{\tiny{[GM$|$GM]}}
                [1, 3, 0]
            }
        }
    }
    & & \\ \\

    \theutterance \stepcounter{utterance}  
    & & & \multicolumn{2}{p{0.3\linewidth}}{
        \cellcolor[rgb]{0.9,0.9,0.9}{
            \makecell[{{p{\linewidth}}}]{
                \texttt{\tiny{[GM$|$GM]}}
                [[1, 3, 0], [0, 0, 1]]
            }
        }
    }
    & & \\ \\

\end{supertabular}
}

\end{document}
